% GENERATED FILE. Edit canonical lessons, then run npm run generate:books.
% canonical-source-hash: fnv1a64:462fd680dc73bc9e
% canonical-lessons: PA-W13-form-mixed-cues, PA-W13-form-timed, PA-W13-message-pacing, PA-W13-message-timed, PA-W13-timed-repair

\chapter{Try the A1 Writing Tasks with Short Clocks}
\label{ch:pa-a1-writing-short-clocks}
\hlchaptermodality{157}

\begin{chapteropening}
\textbf{By the end of this chapter:} \emph{I can complete a six-field Gurmukhi form and a 30–40-word named-reader message in separate short timed checkpoints, then repair content and script in bounded passes without mistaking these drills for a complete A1 mock.}

After a timed six-field mixed-cue form and a short message pacing rehearsal, the learner independently writes a 30–40-word Gurmukhi message for a named reader under a short clock; the full twenty-minute paper remains for the mocks.
\end{chapteropening}

\section[\texorpdfstring{\pa{ਨਾਂ}: · \pa{ਭਾਸ਼ਾ} … (mixed familiar …)}{ਨਾਂ: · ਭਾਸ਼ਾ … (mixed familiar …)}]{Three lines from two old cards}

\label{lesson:PA-W13-form-mixed-cues}



\begin{quote}
The old fictional A · \pa{ਕ} and B · \pa{ਖ} cards are familiar. A form can choose a different old card for each line. The mark tells you which known value to retrieve; it is not a new word or a new person to memorise.
\end{quote}

\subsection*{Writing — supported mixed selection}
Reopen the two old value banks for this rehearsal only. Read one label and its mark, select the taught value, then cover that line before moving to the next. There is no clock and no independent writing credit in this supported step.

\begin{quote}
\pa{ਨਾਂ} (A · \pa{ਕ}): \_\_\_\_\_\_
\end{quote}

\begin{quote}
\pa{ਭਾਸ਼ਾ} (B · \pa{ਖ}): \_\_\_\_\_\_
\end{quote}

\begin{quote}
\pa{ਰਿਹਾਇਸ਼} (A · \pa{ਕ}): \_\_\_\_\_\_
\end{quote}

\subsection*{Your turn — one check}
Check only card selection first. Then check spelling and the space after each colon in a separate pass. The next lesson will close the banks and use all six familiar labels.

\begin{tcolorbox}[breakable,colback=teal!4,colframe=teal!35!black,title={Before you move on}]
Close both banks after the supported check.
\end{tcolorbox}
\section[\texorpdfstring{\pa{ਨਾਂ}: · \pa{ਭਾਸ਼ਾ} … (short timed …)}{ਨਾਂ: · ਭਾਸ਼ਾ … (short timed …)}]{A short clock for six known fields}

\label{lesson:PA-W13-form-timed}



\begin{quote}
Close the old cards. Keep a blank page and a three-minute timer. The six marks below mix familiar values; no model is shown.
\end{quote}

\subsection*{Writing — independent timed form}
Start the three-minute timer. Retrieve each marked value and write all six fields in Gurmukhi, with Gurmukhi digits for age and phone. Stop on time and leave unfinished lines as evidence. Keep models and aids closed.

\begin{quote}
\pa{ਨਾਂ} (A · \pa{ਕ}): \_\_\_\_\_\_
\end{quote}

\begin{quote}
\pa{ਭਾਸ਼ਾ} (B · \pa{ਖ}): \_\_\_\_\_\_
\end{quote}

\begin{quote}
\pa{ਰਿਹਾਇਸ਼} (A · \pa{ਕ}): \_\_\_\_\_\_
\end{quote}

\begin{quote}
\pa{ਕੰਮ} (B · \pa{ਖ}): \_\_\_\_\_\_
\end{quote}

\begin{quote}
\pa{ਉਮਰ} (A · \pa{ਕ}): \_\_\_\_\_\_
\end{quote}

\begin{quote}
\pa{ਫ਼ੋਨ} (B · \pa{ਖ}): \_\_\_\_\_\_
\end{quote}

\subsection*{Your turn — record, then check}
Record completed lines and elapsed seconds. After the timer, check values, spelling, spacing, digit order, and phone grouping separately. Preserve the first attempt; repair is not timed evidence.

\begin{tcolorbox}[breakable,colback=teal!4,colframe=teal!35!black,title={Before you move on}]
This short checkpoint is not the full A1 paper. Two complete mocks test the twenty-minute form-plus-message condition.
\end{tcolorbox}
\section[(named reader, short clock)]{Pace the first two message chunks}

\label{lesson:PA-W13-message-pacing}



\begin{quote}
The whole A1 writing paper lasts twenty minutes. Its message is 30–40 words for one named reader and purpose. Today you will time only a small opening, not pretend to fit that whole paper into a short lesson.
\end{quote}

\subsection*{Writing — one small timed opening}
Close the earlier message pages. Imagine Aman writing to Manan. In ninety seconds, make two familiar chunks: a greeting that names Manan, then one wellbeing question. You may choose the exact taught words and punctuation; no sample sentence is visible here. Stop at ninety seconds and record the number of words you produced.

\subsection*{Your turn — inspect one boundary}
After the timer, check only whether Manan is clearly the reader. Next check the question mark and spacing. Repairing after the clock helps learning but does not change the recorded first attempt.

\begin{tcolorbox}[breakable,colback=teal!4,colframe=teal!35!black,title={Before you move on}]
The next checkpoint uses a fresh purpose, a full 30–40-word message, and a short clock. The old supported message lacks the required water detail, so copying its old order would miss that purpose.
\end{tcolorbox}
\section[(timed independent named-reader message)]{A new message under a short clock}

\label{lesson:PA-W13-message-timed}



\begin{quote}
Close all earlier message pages. Set a three-minute timer and keep a blank page. The task below uses known language, but the reader, speaker, and purpose combination is not the old independent note. No romanization, sentence bank, or copyable answer is available while the clock runs.
\end{quote}

\subsection*{Writing — independent timed message}
As Aman, write Manan a \textbf{30–40-word} Gurmukhi message. Greet Manan and ask about his wellbeing. Say that you like water and one other familiar thing. Then suggest meeting again and close in a fitting tone. You may add other already taught words to reach the band. Choose and order your own sentences; the old supported Aman-to-Manan model did not include the water purpose.

Start the timer. Stop at three minutes, even if unfinished. Record the word count and elapsed time on a separate line. Do not use a dictionary, translator, spell-checker, or earlier page during this first attempt.

\subsection*{Your turn — record before repair}
Before reopening anything, note whether the named reader, water detail, meeting purpose, and 30–40-word band are present. Keep the first attempt unchanged. Repair comes in the next short lesson, not inside the timed score.

\begin{tcolorbox}[breakable,colback=teal!4,colframe=teal!35!black,title={Before you move on}]
This short message checkpoint is timed writing evidence, not a complete twenty-minute A1 paper. The two full mocks and human validation remain separate gates to readiness.
\end{tcolorbox}
\section[(repair after the short clocks)]{Repair without rewriting the clock}

\label{lesson:PA-W13-timed-repair}



\begin{quote}
Keep the original timed form and message unchanged. Copy neither into a new scored answer. The clock showed what you could produce then; this short lesson shows what one careful repair can teach next.
\end{quote}

\subsection*{Your turn — separate checks}
First inspect task fulfilment: six requested form values, one named reader, the requested message details, and a 30–40-word message. Then inspect order: each value beside its own label, and greeting, content, invitation, and close in a readable sequence. Third, inspect Punjabi vocabulary and grammar, including agreement and tone. Last, inspect Gurmukhi spelling, word spaces, punctuation, digit order, and phone grouping. Do one pass at a time.

Choose one error dimension and write one corrected line on a separate page. Say what changed. Do not present this later line as the timed response.

\begin{tcolorbox}[breakable,colback=teal!4,colframe=teal!35!black,title={Before you move on}]
Record one next practice target for the form and one for the message. A full paper still asks for both tasks in twenty minutes; only the complete timed mocks can test that combined condition.
\end{tcolorbox}
